% ============================================================
% Project: THE PHILOSOPHY OF VIABILITY
% Book: Book I — Ontology of Viability
% Book ID: VS.B01
% Part: Part I — Identity and Manifestation
% Chapter: Chapter 07 — Formal and Logical Reconstruction
% File: 2_MAINMATTER/PART_I_IDENTITY_AND_MANIFESTATION/CHAPTER_07_FORMAL_AND_LOGICAL_RECONSTRUCTION/sections/section_02_ontological_vocabulary_and_primitive_terms.tex
% Section: 02 — Ontological Vocabulary and Primitive Terms
% Purpose: To establish and classify the principal ontological
% terms required for formal reconstruction, distinguish conceptual
% primitives from formal primitives, and specify the conditions
% governing their legitimate formal use.
% ============================================================

\section{Ontological Vocabulary and Primitive Terms}
\label{sec:ch07_s02}

\subsection{The Need for a Disciplined Ontological Vocabulary}

Formal reconstruction cannot proceed adequately without a
disciplined account of the vocabulary through which its subject
matter is represented. The choice of terms is not a merely
terminological matter. It determines which distinctions can be
expressed, which relations can be examined, and which inferences
may appear available within the formal framework.

A formal language may be internally consistent while remaining
philosophically inadequate if its vocabulary conflates concepts
that the ontology requires to distinguish. Conversely, a
vocabulary may contain many carefully differentiated terms
without yet providing a coherent account of how those terms
relate to one another. The task is therefore not to maximize
the number of technical expressions, but to establish a
controlled vocabulary whose terms have identifiable functions
and whose relations are not assumed without justification.

The preceding chapters have established that identity must be
examined in relation to being, determination, distinction,
unity, plurality, invariance, expression, manifestation,
persistence, and continuity. These concepts belong to a
connected ontological inquiry, but their relevance to identity
does not make them interchangeable. Their formal treatment
requires that their respective roles remain explicit.

The present section establishes the vocabulary needed for
that treatment. It distinguishes three matters that must not
be confused:

\begin{enumerate}
    \item the ontological concepts and phenomena that constitute
    the subject matter of the inquiry;
    \item the conceptual terms through which those phenomena
    are analytically articulated;
    \item the formal symbols and expressions through which
    selected claims about them are represented.
\end{enumerate}

These matters are related, but they do not occupy the same
level of analysis. An ontological concept is not identical
to the word used to express it, and neither the word nor
its formal symbol is identical to the ontological subject
matter to which it refers.

The vocabulary developed here must consequently be adequate
to the ontology without being mistaken for the ontology itself.

\subsection{Ontological Terms, Conceptual Terms, and Formal Symbols}

An ontological inquiry concerns what exists and the conditions,
relations, and determinations belonging to what exists. Its
terms are used to articulate aspects of this subject matter.
A formal reconstruction subsequently represents selected
claims through a language whose expressions are governed by
specified rules.

Three levels must therefore be distinguished.

At the ontological level, the inquiry concerns being, existents,
identity, determination, and the relations or conditions
attributed to them. These are the matters about which the
theory makes claims.

At the conceptual level, these matters are distinguished,
compared, and articulated through concepts. The terms
\emph{being}, \emph{identity}, \emph{determination},
\emph{unity}, and \emph{manifestation}, for example, are
used to examine different aspects of the ontological subject
matter. Their conceptual distinction is necessary for
analysis, but does not establish that the corresponding
ontological dimensions exist as independently separable
entities.

At the formal level, symbols, variables, predicates, relations,
operators, and rules are introduced to represent selected
conceptual claims. Their meaning depends upon the
interpretation assigned to them and the formal rules
governing their use.

The distinction may be expressed schematically:

\[
\text{Ontological Subject Matter}
\longrightarrow
\text{Conceptual Articulation}
\longrightarrow
\text{Formal Representation}
\]

The arrows indicate methodological articulation and
representation. They do not indicate that the ontological
subject matter is produced by concepts or constituted by
formal symbols. Nor do they imply that every ontological
phenomenon can be exhaustively represented within one
formal language.

% ============================================================
% Project: THE PHILOSOPHY OF VIABILITY
% Book: Book I — Ontology of Viability
% Book ID: VS.B01
% Part: Part I — Identity and Manifestation
% Chapter: Chapter 07 — Formal and Logical Reconstruction
% File: 4_ASSETS/diagrams/VS.B01.P01.C07.S02.SS02.P001.DG01.levels_of_ontological_and_formal_articulation.tex
% Asset: DG01 — Levels of Ontological and Formal Articulation
% Purpose: To distinguish ontological subject matter,
% conceptual articulation, and formal representation.
%
% Requires (preamble):
%   \usepackage{tikz}
%   \usetikzlibrary{arrows.meta, positioning}
% ============================================================

\begin{figure}[htbp]
    \centering
    \hyphenpenalty=10000 \exhyphenpenalty=10000\relax  % no hyphenation inside diagram labels

    % separator used inside the level boxes
    \def\vsdot{\enspace$\cdot$\enspace}

    \begin{tikzpicture}[
        font=\small,
        node distance=1.6cm,
        % --- base box ---
        box/.style={
            draw, rounded corners=3pt, align=center,
            text width=8.0cm, minimum height=1.9cm,
            inner sep=7pt, line width=0.8pt
        },
        % --- node roles ---
        ontbox/.style={box, fill=blue!8,  draw=blue!55!black},
        concbox/.style={box, fill=blue!16, draw=blue!55!black},
        formbox/.style={box, fill=gray!10, draw=black!60},
        % --- arrows (methodological relations, not production) ---
        directed/.style={-{Stealth[length=2.5mm]}, line width=0.9pt},
        feedback/.style={-{Stealth[length=2.5mm]}, line width=0.9pt, dashed},
        elabel/.style={font=\footnotesize\itshape, text=black!75,
                       align=left, text width=3.0cm}
    ]

    % ---------------------------------------------------------------
    % Three levels
    % ---------------------------------------------------------------
    \node[ontbox] (ontological) {
        \textbf{Ontological Subject Matter}\\[3pt]
        Being\vsdot Existent\vsdot Identity\vsdot Determination\\
        Relations, conditions, and manifestations
    };
    \node[concbox, below=of ontological] (conceptual) {
        \textbf{Conceptual Articulation}\\[3pt]
        Terms\vsdot Distinctions\vsdot Definitions\vsdot Claims\\
        Interpretations and ontological arguments
    };
    \node[formbox, below=of conceptual] (formal) {
        \textbf{Formal Representation}\\[3pt]
        Variables\vsdot Predicates\vsdot Relations\vsdot Operators\\
        Rules, models, and inferential structures
    };

    % ---------------------------------------------------------------
    % Articulation (solid, downward, left) and critical feedback
    % (dashed, upward, centre-right); labels sit beside the arrows
    % ---------------------------------------------------------------
    \draw[directed] ([xshift=-3.3cm]ontological.south) --
        node[elabel, right=5pt] {is articulated as a subject of inquiry}
        ([xshift=-3.3cm]conceptual.north);
    \draw[directed] ([xshift=-3.3cm]conceptual.south) --
        node[elabel, right=5pt] {is represented through formal expressions}
        ([xshift=-3.3cm]formal.north);

    \draw[feedback] ([xshift=0.8cm]conceptual.north) --
        node[elabel, right=5pt] {ontological adequacy constrains conceptual articulation}
        ([xshift=0.8cm]ontological.south);
    \draw[feedback] ([xshift=0.8cm]formal.north) --
        node[elabel, right=5pt] {formal analysis may require conceptual revision}
        ([xshift=0.8cm]conceptual.south);

    % ---------------------------------------------------------------
    % Legend
    % ---------------------------------------------------------------
    \node[below=0.6cm of formal, font=\scriptsize] {
        \tikz[baseline=-0.5ex]\draw[directed] (0,0) -- (0.7,0);~articulation
        \quad\quad
        \tikz[baseline=-0.5ex]\draw[feedback] (0,0) -- (0.7,0);~critical feedback
    };

    \end{tikzpicture}

    \caption{The distinction between ontological subject matter, conceptual articulation, and formal representation. The arrows indicate methodological relations and critical feedback, not the production of ontological reality by concepts or formal systems.}
    \label{fig:ch07_levels_of_articulation}
\end{figure}

% ============================================================
% END OF ASSET DG01 — CHAPTER 07, SECTION 02
% ============================================================


The diagram distinguishes the subject matter of the inquiry
from the conceptual and formal means through which claims
about it are articulated. The distinction is essential to
an ontology that does not make the existence or identity
of an existent dependent upon its recognition, description,
or representation.

A formal expression therefore receives its significance
through an interpretation that connects it to the concepts
under examination. The interpretation must be stated
rather than silently assumed. If a symbol is used to
represent identity, for example, the reconstruction must
clarify whether it denotes identity as an ontological
dimension, a relation of numerical identity, a criterion
of identification, or a particular formal relation.
The mere use of a symbol cannot settle that question.

\subsection{Conceptual Primitives and Formal Primitives}

The expression \emph{primitive term} can be used in more
than one sense. Unless these senses are distinguished,
the formal reconstruction risks turning a methodological
decision into an unsupported ontological thesis.

A \emph{conceptual primitive} is a term whose meaning,
within the present framework, is not established by
reducing it to a more basic concept already available
in that framework. Its role is sufficiently fundamental
that the inquiry must employ it in order to articulate
other concepts and claims. A conceptual primitive may
still be clarified through explanation, contrast,
examples, and examination of its relations. Its
primitive status does not make it exempt from philosophical
analysis.

A \emph{formal primitive}, by contrast, is an expression
that is not defined through other expressions within
a particular formal language. It is assigned an
interpretation or role and then used according to the
rules of that language. Its status as formally primitive
is relative to the chosen system. A different formal
system may define the same expression through other
terms or represent the relevant subject matter in
another way.

These two senses must not be identified.

If a term is left formally undefined, this does not
establish that its ontological referent is absolutely
simple, indivisible, or incapable of further philosophical
analysis. It establishes only that the selected formal
language does not define it through other expressions.
Likewise, a concept may be philosophically fundamental
without being represented by a single uninterpreted
symbol in every formal reconstruction.

The present chapter therefore adopts a restrained
approach to primitive vocabulary. It identifies a
limited set of terms that must be available for the
formal examination of the ontology, while leaving
their philosophical interpretation open to the
arguments already developed and those still required.

In particular, the terms \emph{being}, \emph{existent},
and \emph{identity} occupy foundational positions
in the present inquiry. Their foundational role
does not authorize their collapse into one another.
Being concerns the actuality of what exists; identity
concerns the existent as the existent that it is;
and the term \emph{existent} refers to that which
is under ontological consideration. These formulations
establish their respective analytical roles without
treating them as three independently existing
components.

The relation between being and identity remains
governed by the three distinct commitments established
in Chapter 01:

\begin{enumerate}
    \item \textbf{Analytical Distinction:} being and
    identity are distinguishable in analysis.

    \item \textbf{Ontological Inseparability:} being
    and identity are not treated as independently
    existing actualities detachable from one another.

    \item \textbf{Absence of Ontological Priority:}
    neither being nor identity is assigned ontological
    priority over the other.
\end{enumerate}

These commitments must not be compressed into one
undifferentiated formula. Each expresses a different
aspect of the relation and requires its own justification.

The vocabulary of the reconstruction must preserve
this independence of commitments. A formal representation
that expresses inseparability but cannot distinguish
the analytical roles of being and identity would be
inadequate. A representation that distinguishes them
but introduces a hierarchy of ontological dependence
would likewise fail to preserve the adopted position.

\subsection{The Foundational Ontological Vocabulary}

The foundational vocabulary consists of terms required
to articulate the subject matter of Part I and to
formulate the principal questions examined in this
chapter. Their inclusion does not imply that each
has already received a complete definition.

\paragraph{Being.}

\emph{Being} refers to the actuality of what exists.
It is used to articulate the ontological actuality
of an existent. In the present framework, being is
analytically distinguishable from identity but
ontologically inseparable from it. The term must
not be used as though it denoted an identity-neutral
substrate that is first constituted and subsequently
receives identity.

\paragraph{Existent.}

An \emph{existent} is that which is considered as
existing and as being the existent that it is.
The term identifies the subject of ontological
attribution and analysis. It does not, by itself,
settle whether the existent should be interpreted
as a substance, aggregate, process, relation,
composition, or another ontological configuration.
Such classifications require further argument.

\paragraph{Identity.}

\emph{Identity} concerns the existent as the existent
that it is. It is a foundational ontological concern
of this work and cannot be reduced to an additional
property attributed to an already constituted existent.
Nor can it be identified without argument with
determination, unity, singularity, individuation,
numerical identity, invariance, structure, composition,
signature, manifestation, persistence, continuity,
or stability.

The term must remain available for the ontological
question itself. A formal representation may clarify
relations involving identity, but it must not assume
that one of those relations exhausts the meaning
or constitution of identity.

\paragraph{Determination.}

\emph{Determination} concerns the determinate
character of an existent. It addresses the ways
in which an existent is ontologically specified
rather than wholly unspecified. Determination is
indispensable to the examination of identity, but
identity is not identical with the totality of
determinations associated with an existent.

The distinction between determination in general
and identity-constitutive determination must be
preserved. The presence, persistence, visibility,
or descriptive usefulness of a determination does
not automatically establish its constitutive role
in identity. That role requires an independent
ontological argument.

\paragraph{Distinction and difference.}

\emph{Distinction} concerns the analytical or
ontological differentiation of terms, aspects,
or existents, according to the respect under
consideration. \emph{Difference} concerns a
specified manner in which terms or existents
differ. Neither term can be used as an automatic
substitute for identity or identity distinction.

A formal reconstruction must specify whether a
distinction is conceptual, logical, relational,
or intended as an ontological distinction. It
must not infer ontological separation merely
from the fact that two concepts can be distinguished
in analysis.

\paragraph{Relation.}

A \emph{relation} connects or orders terms according
to a specified respect. Its formal representation
requires the identification of the relata, the
relevant domain, and the intended relational role.
The existence of a relation does not, by itself,
establish whether the relation is constitutive
of identity, dependent upon already distinguishable
identities, or relevant in another way.

The term must therefore be used with explicit
attention to its arguments and interpretation.
An unspecified relation symbol cannot carry an
ontological conclusion that has not been established.

\paragraph{Manifestation.}

\emph{Manifestation} concerns the manifest presence
or presentation of an existent or of aspects
associated with it. Its relation to identity
requires careful examination. A manifestation
may disclose or articulate aspects of an existent,
but the formal description of a manifestation
does not, by itself, exhaust the identity of
that existent.

\paragraph{Expression.}

\emph{Expression} concerns the articulation or
expression of aspects associated with an existent.
It must be distinguished from manifestation,
even where the two are closely related in a
particular ontological account. The formal
reconstruction must not use one term as a
synonym for the other without specifying the
relation intended.

These foundational terms provide the vocabulary
required for the subsequent formal analysis. Their
roles are summarized in Table~\ref{tab:ch07_ontological_vocabulary}.

% ============================================================
% Project: THE PHILOSOPHY OF VIABILITY
% Book: Book I — Ontology of Viability
% Book ID: VS.B01
% Part: Part I — Identity and Manifestation
% Chapter: Chapter 07 — Formal and Logical Reconstruction
% File: 4_ASSETS/tables/VS.B01.P01.C07.S02.SS04.P001.TL01.ontological_vocabulary_and_formal_roles.tex
% Asset: TL01 — Ontological Vocabulary and Formal Roles
% Purpose: To distinguish foundational ontological terms,
% related concepts, and formal constraints governing their use.
% ============================================================

\begin{table}[htbp]
\centering
\small
\caption{Ontological Vocabulary and Its Formal Roles}
\label{tab:ch07_ontological_vocabulary}
\renewcommand{\arraystretch}{1.25}
\begin{tabular}{|p{0.19\textwidth}|p{0.29\textwidth}|p{0.42\textwidth}|}
\hline
\textbf{Term} &
\textbf{Role in the Ontological Inquiry} &
\textbf{Constraint on Formal Use}
\\
\hline
Being &
The actuality of what exists. &
Must remain analytically distinguishable from identity while ontologically inseparable from it.
\\
\hline
Existent &
That which is considered as existing and as the existent that it is. &
Must not be presupposed to be exclusively a substance, aggregate, process, or composition.
\\
\hline
Identity &
The existent as the existent that it is. &
Must not be reduced to determination, unity, numerical identity, invariance, manifestation, or continuity.
\\
\hline
Determination &
The determinate character of an existent. &
The constitutive role of a particular determination must be established rather than inferred from its presence.
\\
\hline
Distinction and Difference &
Analytical or ontological differentiation under a specified respect. &
Conceptual distinction must not be treated as proof of ontological separation; difference must be specified.
\\
\hline
Relation &
A connection or ordering among terms under a specified respect. &
The relata, domain, and intended relational role must be identified.
\\
\hline
Manifestation &
The manifest presence or presentation of an existent or its aspects. &
Must not be treated as an exhaustive substitute for the identity of the existent.
\\
\hline
Expression &
The articulation or expression of aspects associated with an existent. &
Must remain distinguishable from manifestation unless their relation is explicitly established.
\\
\hline
Unity and Plurality &
Multiplicity or unity considered under a specified respect. &
Component, manifestational, relational, and identity plurality must not be conflated.
\\
\hline
Invariance &
Preservation under a specified transformation and respect. &
Transformation, domain, respect, type, and limits must be specified.
\\
\hline
Transformation &
Change between conditions or configurations. &
Must not be taken to establish either preservation or replacement of identity by itself.
\\
\hline
Persistence and Continuity &
Continuation across an interval or sequence of conditions. &
Must be distinguished from identity and specified with respect to subject, interval, and conditions.
\\
\hline
Structure and Composition &
Organization of relations or constitution through components. &
Neither structure nor composition may be assumed to exhaust identity.
\\
\hline
Numerical Identity &
Being one and the same object within a specified logical interpretation. &
The logical equality symbol must not be treated as a complete expression of ontological identity.
\\
\hline
Signature &
A structured characterization associated with an existent. &
Descriptive or identifying usefulness does not establish ontological constitution.
\\
\hline
\end{tabular}
\end{table}

% ============================================================
% END OF ASSET TL01 — CHAPTER 07, SECTION 02
% ============================================================

The table is a controlled vocabulary for the present
reconstruction. It does not claim that every term
has been exhaustively defined, nor that the listed
terms constitute a complete inventory of ontology.
Its function is to establish the minimum distinctions
that must be preserved when claims are translated
into formal expressions.

\subsection{Related Concepts and Their Distinct Functions}

Several concepts are indispensable to the inquiry
because they concern the conditions, organization,
or possible transformations of existents. They
must nevertheless remain distinct from the
foundational vocabulary.

\paragraph{Unity.}

\emph{Unity} concerns the way in which an existent
or configuration is considered as a unity. It
does not automatically imply singularity, absence
of internal differentiation, or the existence
of one and only one identity. Unity must not
be equated with aggregate or composition.

\paragraph{Plurality.}

\emph{Plurality} concerns multiplicity under a
specified respect. The relevant respect may
involve components, manifestations, relations,
determinations, or identities. A plurality
established in one respect cannot be transferred
to another without an additional argument.
In particular, plurality of components or
manifestations does not by itself establish
plurality of identities.

\paragraph{Singularity.}

\emph{Singularity} concerns a condition of being
one or singular under a specified interpretation.
It must not be treated as synonymous with unity
or identity. A claim of singularity requires
clarification of what is being counted and under
which respect.

\paragraph{Individuation.}

\emph{Individuation} concerns the determination
or differentiation of an existent as this existent
rather than another. It is relevant to the
question of identity, but it cannot be substituted
for identity itself. The formal reconstruction
must distinguish the conditions under which
individuation is discussed from the ontological
status of identity.

\paragraph{Numerical identity.}

\emph{Numerical identity} concerns whether two
designations, occurrences, or references concern
one and the same existent. It is especially
important in formal logic, where identity
expressions may be governed by established
logical rules. However, numerical identity
does not automatically provide a complete
account of the ontological constitution of
identity developed in this work.

\paragraph{Invariance.}

\emph{Invariance} concerns what remains unchanged
under a specified transformation, within a
specified domain, and with respect to a specified
respect or type. Every formal claim of invariance
must make these conditions explicit. Formal
invariance does not by itself establish ontological
invariance, and neither is equivalent to identity.

\paragraph{Transformation.}

\emph{Transformation} concerns a change from one
condition, configuration, or state to another
under specified conditions. A transformation
must be described with sufficient precision to
identify what changes, what is compared, and
which domain is involved. The occurrence of
transformation does not independently establish
either the preservation or replacement of identity.

\paragraph{Persistence and continuity.}

\emph{Persistence} concerns the continued existence
of an existent across a relevant interval or
sequence of conditions. \emph{Continuity} concerns
a relation of connectedness or continuation
across the relevant stages or conditions.
Neither term is synonymous with identity.
A claim about persistence or continuity must
specify the subject, interval, conditions, and
respect under which it is made.

\paragraph{Structure and composition.}

\emph{Structure} concerns an organization of
elements or relations. \emph{Composition} concerns
the constitution of a configuration through
components or constituents. Neither concept
can be assumed to exhaust identity. A structural
or compositional account may be relevant to
identity, but its constitutive significance
must be established.

\paragraph{Signature.}

\emph{Signature} concerns a structured
characterization associated with an existent
or configuration. It may be useful for
identification, comparison, or formal description.
Its usefulness does not establish that it is
identical with, or sufficient to constitute,
the identity of the existent.

The distinctions among these concepts are not
introduced to isolate them from identity. They
are introduced to prevent an unjustified reduction
of identity to any one of the phenomena through
which it may be examined.

\subsection{Terms Requiring Explicit Relational Specification}

Some terms do not become sufficiently determinate
for formal use until the relations in which they
participate have been specified. Their use without
such specification risks producing expressions
that appear precise but remain semantically
underdetermined.

A claim concerning invariance, for example, is
incomplete if it does not identify the transformation
under consideration. It must also specify the
domain over which the claim is made, the respect
in which invariance is asserted, the type of
feature or relation that remains invariant,
and the limits within which the claim is intended
to hold.

Similarly, a claim concerning identity plurality
must specify the existent or configuration
under consideration, the respect in which
plurality is asserted, and the grounds for
distinguishing the identities involved. A count
of components, manifestations, or relations
cannot replace these specifications.

A claim concerning continuity must identify
what is said to continue, across which interval
or sequence of conditions, and under which
criterion of connectedness. It must not silently
assume that continuity establishes identity.

A claim concerning manifestation must identify
what is manifested, the relevant conditions
of manifestation, and the relation between
the manifestation and the existent or identity
under examination. The manifestation must
not be treated as a complete substitute for
the manifested existent.

These requirements indicate that formal
precision depends not merely on the selection
of terms but on the specification of their
relational roles. A symbol without a defined
argument structure or interpretation is not
made adequate by its symbolic appearance.

\subsection{The Distinction Between Ontological Identity and Logical Equality}

A particularly important terminological issue
arises from the use of the word \emph{identity}
in formal logic.

In classical first-order logic, the identity
symbol, commonly written as \(=\), expresses
numerical identity between the objects assigned
to its terms. It is governed by formal principles
that specify how identical objects may be
substituted within appropriate contexts.

This logical use is indispensable for formal
reasoning, but it must not be assumed to exhaust
the ontological problem of identity addressed
in Part I.

The question whether two expressions designate
one and the same object is not identical to
the question of what it means for an existent
to be the existent that it is. The former
concerns numerical identity within a logical
language and its interpretation. The latter
concerns the ontological status and constitution
of identity.

The distinction can be illustrated by considering
two different questions:

\begin{enumerate}
    \item Do the expressions \(a\) and \(b\)
    designate one and the same object within
    a specified interpretation?

    \item What is the ontological significance
    of the identity of the existent designated
    by those expressions?
\end{enumerate}

The first question can be represented through
the logical identity relation. The second
requires an ontological account that cannot
be obtained merely by applying the formal
rules governing \(=\).

Accordingly, the equality symbol will be used
in this chapter only where numerical identity
within a specified formal language is intended.
It will not be used as an unexplained symbol
for the full ontological concept of identity.

Where the reconstruction requires reference
to identity as an ontological dimension, the
notation and interpretation must be introduced
explicitly. No symbolic convention may silently
collapse the distinction between ontological
identity and logical equality.

This restriction is not a rejection of classical
logic. It is a requirement that the formal
resources of logic be used with due attention
to the philosophical question they are intended
to represent.

\subsection{Definition, Explication, and Theoretical Commitment}

The vocabulary of a formal reconstruction
requires definitions, but not every important
term can be adequately handled through a
stipulative definition.

A stipulative definition establishes how an
expression is to be used within a specified
discussion or formal system. It may be useful
for introducing notation, delimiting a technical
term, or avoiding ambiguity. Its validity
as a stipulation depends upon clarity and
consistency of use.

An explication attempts to clarify a concept
by making its intended role, conditions of
application, and relations to other concepts
more precise. Explication may involve revision
of ordinary usage, but it must remain answerable
to the philosophical problem that motivates it.

A theoretical commitment, by contrast, asserts
or adopts a claim concerning the subject matter.
It is not established merely by defining the
terms in which it is expressed.

These functions must remain distinct. If the
expression \emph{identity-constitutive determination}
is introduced as a technical term, the
introduction of that term does not establish
which determinations actually satisfy the
relevant condition. If a formal relation is
defined to represent a particular connection,
the definition does not establish that the
connection obtains in reality. If a term is
stipulated to designate a class of configurations,
the stipulation does not prove that the
corresponding configurations exist.

The reconstruction must therefore distinguish
the introduction of terminology from the
justification of ontological claims.

A definition may be evaluated for clarity,
consistency, scope, and usefulness. An
ontological proposition requires assessment
of its grounds and implications. The two
evaluations may interact, but they must not
be conflated.

\subsection{Vocabulary Discipline and Conditions of Use}

The vocabulary established in this section
is governed by several requirements.

First, every term must retain a sufficiently
stable meaning within the argument in which
it is used. Where a term has more than one
legitimate meaning, the relevant sense must
be identified.

Second, the same formal symbol must not be
used for different ontological roles unless
the distinction is explicitly represented
in its interpretation. Notational economy
cannot justify conceptual conflation.

Third, the introduction of a new technical
term must serve a specific analytical purpose.
A new term is warranted when it clarifies
a distinction, identifies a previously
unarticulated relation, or makes an argument
more precise. It is not warranted merely
because an existing term has been assigned
a new label.

Fourth, a term must not acquire a stronger
ontological meaning merely through repeated
use. Repetition does not establish necessity,
universality, or constitutive status.

Fifth, relations among terms must be stated
with appropriate scope. A relation established
under one respect cannot be generalized to
every respect without argument.

Sixth, the formal reconstruction must preserve
the distinction between an ontological claim
and a claim about the representation of that
claim. The formal properties of an expression
belong to the language and system in which
it occurs; the ontological significance
of the expression depends upon its justified
interpretation.

Seventh, unresolved terms and relations must
remain visibly unresolved. Where the preceding
chapters have established a question but not
a definitive answer, formal reconstruction
must not conceal that status by introducing
an apparently exact definition.

These requirements establish a vocabulary
that is disciplined without being prematurely
closed. They permit formal analysis while
preserving the possibility that further
philosophical inquiry may require the
refinement, extension, or revision of
particular terms.

\subsection{The Limits of the Present Primitive Vocabulary}

The vocabulary identified in this section
is sufficient to begin the formal reconstruction
of the principal claims of Part I. It is
not presented as a final inventory of all
ontological primitives.

Several questions remain open.

It has not yet been established whether
every ontological relation can be represented
through a single general relational form,
or whether different kinds of relations
require distinct formal types.

It has not yet been established whether
determination should be represented as a
property, a relation, a structured configuration,
or through more than one formal category.
The choice must depend upon the ontological
role being represented and must not be made
merely for notational convenience.

It has not yet been established whether
the concept of identity can be represented
through a single formal predicate, through
a family of related expressions, or through
a more complex formal structure. Any such
representation must preserve the distinction
between identity and its related concepts.

It has not yet been established whether
the relations among being, identity, and
determination can be adequately expressed
within one formal language without loss
of their respective conceptual roles.

Nor has it been established that the
vocabulary of classical first-order logic
is sufficient for every question raised
by the ontology. The adequacy of a formal
language must be assessed in relation to
the claims it is intended to represent.

These matters are not defects to be concealed.
They are methodological questions that
must remain open until the subsequent
analysis supplies sufficient grounds for
a more determinate formal treatment.

\subsection{Conclusion: A Controlled Vocabulary for Formal Reconstruction}

The present section has established the
need for a controlled ontological vocabulary
and distinguished the ontological subject
matter from its conceptual articulation
and formal representation.

It has identified being, existent, identity,
determination, distinction, relation,
manifestation, and expression as foundational
terms for the present reconstruction, while
preserving the distinct functions of unity,
plurality, singularity, individuation,
numerical identity, invariance, transformation,
persistence, continuity, structure, composition,
and signature.

It has also distinguished conceptual primitives
from formal primitives. The primitive status
of a term within a formal language does not
establish the absolute simplicity or
ontological irreducibility of its referent.
Likewise, the formal identity symbol must
not be mistaken for a complete representation
of ontological identity.

The resulting vocabulary is deliberately
controlled rather than exhaustive. Its
purpose is to make subsequent formal
representations more precise while
preventing the formal language from
introducing commitments that the ontology
has not established.

The next section examines the formal domains
and types through which this vocabulary
may be represented. It will consider how
the formal language can distinguish objects,
properties, relations, transformations,
and other relevant expressions without
assuming that the categories of the formal
system automatically correspond to
independent ontological divisions.

% ============================================================
% END OF SECTION 02 OF CHAPTER 07
% ============================================================